\section{Introduction}

Payment stablecoins are money-like liabilities backed by portfolios that look increasingly like those of
government money market funds. The two largest issuers alone report well over \$200\,bn of US Treasury bills,
reverse repo and government money market fund shares, and the GENIUS Act signed in July 2025 confines permitted
reserve assets to cash, insured deposits, bills with no more than 93 days to maturity, Treasury-backed repo and
government money market funds \citep[][all institutional figures used in this paper are listed for verification in
the repository]{genius2025}. A growing empirical literature documents that stablecoin flows already move the front
end of the Treasury curve in normal times \citep{ahmed2025}, and the Treasury Borrowing Advisory Committee has
discussed stablecoins as a structural source of bill demand \citep{tbac2025}.

What happens to that demand when a stablecoin runs? Two episodes give partial answers. In May 2022 Tether redeemed
roughly a tenth of its supply in two weeks after the collapse of Terra, with a brief and shallow depeg. In March
2023 USDC fell to about 88 cents over a weekend after Circle disclosed that \$3.3\,bn of its reserves sat at
Silicon Valley Bank, recovered within two days once deposits were guaranteed, and still lost about a quarter of
its supply over the following month \citep{circle2023,tether2022}. Neither episode disturbed the bill market, but
neither involved an issuer of today's size, let alone the multiples that industry projections contemplate. The
question this paper asks is quantitative: when a run happens, how much bill-market dislocation does the forced
liquidation of reserves cause, through which channel, and at what sector size does it become a financial-stability
concern?

Two episodes are too few for econometrics, so we answer with a calibrated agent-based model. Holders are
heterogeneous in size, information, attention and access to the issuer: institutions and arbitrageurs can redeem
at par, retail and most exchanges can only sell on a secondary market. Exit decisions follow a global-games-style
threshold rule on a private signal about reserve impairment, a visible depeg, peer behaviour and a suspension
flag. The secondary price carries a fundamental discount (the expected loss share of whoever remains) and a
liquidity discount proportional to the stock of tokens that arbitrageurs have not yet absorbed. The issuer pays
redemptions through a waterfall (cash, repo, money market fund shares, bills shortest-first, other assets at a
haircut), subject to a daily processing capacity, and manages liquidity by keeping cash for pending claims and
rolling the rest of its maturing bills. The bill market absorbs outright sales and demand withdrawn at auction
with a persistent and a temporary impact component, a dealer-capacity penalty and a reinvestment leak through
which redeemed cash returns via money market funds. A second issuer receives a belief spillover from the first.
The model nests a closed-form threshold aggregate when the behavioural channels are switched off, and we show
that the simulation reproduces it exactly.

Our contribution is threefold. First, we quantify the spillover and identify its channel. At today's size a run on
the largest issuer, whether a liquidity scare or a 10\% loss of reserves, moves three-month bill yields by
17--21\,bp in the model, and the issuer sells no bills outright: with processing capacity of about 1.2\% of
supply per day, redemptions are funded entirely by the \$2.8\,bn a day of maturing bills. The externality is
demand withdrawn at auction, not a fire sale, and the issuer's fire-sale loss is negligible because the bills it
would sell have a duration of weeks. Second, we turn ``stablecoins might be systemic one day'' into a number. Under
square-root impact, the 25\,bp threshold is crossed at 1.4--2 times today's size and the 50\,bp threshold at
5.5--8 times; under linear impact, the upper bound, both are crossed already today. Third, we rank levers. A
minimum cash buffer of 20--30\% removes the footprint of a liquidity run and cuts that of a solvency run by about
two thirds; the 93-day maturity cap converts outright sales into orderly non-rollover; redemption fees do
little; redeeming at net asset value rather than par halves the depeg; a Fed backstop that halves impact halves
the footprint by construction. Sensitivity analysis shows a clean separation: the size of a run is governed by
behavioural parameters (thresholds, attention, the perceived probability of loss), whereas its Treasury-market
footprint is governed by plumbing (how much redeemed cash returns to bills, arbitrage capital, the accounting of
non-rollover) and barely by behaviour at all.

The rest of the paper proceeds as follows. Section~\ref{sec:background} sets out the institutional mechanics.
Section~\ref{sec:lit} relates the paper to the literature. Section~\ref{sec:model} presents the model,
Section~\ref{sec:calib} its calibration and validation, Section~\ref{sec:results} the results,
Section~\ref{sec:robust} robustness, Section~\ref{sec:policy} policy implications. Section~\ref{sec:conclusion}
concludes. Appendices give the parameter table, pseudo-code and additional figures. Code, configuration and every
table in this paper are reproducible from the repository with one command.

\section{Institutional background}\label{sec:background}

\paragraph{Redemption mechanics.} A payment stablecoin is a token redeemable at par from its issuer, but direct
redemption is available only to verified counterparties above a minimum size: for the largest offshore issuer,
verified customers with a minimum of \$100,000 and a 0.1\% fee; for the largest US-regulated issuer, institutional
customers of its minting platform \citep{tether_attest,circle_reserve}. Everyone else obtains dollars by selling the
token on exchanges. The peg is therefore maintained by arbitrageurs who buy discounted tokens on the secondary
market and redeem them at par, and the depth of that arbitrage is bounded by their capital and by the issuer's
processing capacity \citep{ma2023,lyons2023}. During the USDC weekend of March 2023, redemption was unavailable
because banks were closed, arbitrage capital stood aside, and the secondary price cleared at a discount of about
twelve cents. Once redemption resumed and the Silicon Valley Bank exposure was resolved, the discount closed in a
day, but holders who had decided to leave kept redeeming through the following weeks.

\paragraph{Reserve rules.} The GENIUS Act requires one-for-one backing in US coins and currency, insured deposits,
Treasury bills with at most 93 days to maturity, Treasury-backed repo and reverse repo, shares of government money
market funds invested in the same assets, and central bank reserves; issuers must publish monthly reserve reports
and disclose their redemption policy, and may not pay interest to holders \citep{genius2025}. Issuer attestations at
mid-2026 report reserve mixes that already resemble this list: the largest offshore issuer holds roughly two thirds
of its assets in bills directly, with reverse repo, money market fund shares and a residual of other assets
(bitcoin, gold, secured loans); the largest US-regulated issuer holds about an eighth in bank deposits and the rest
in a government money market fund that invests in short bills and overnight repo \citep{tether_attest,circle_reserve}.
These are the two reserve profiles we model as ``Issuer~A'' and ``Issuer~B''; the figures are defaults from public
reports as remembered by the authors and must be replaced with the latest attestation before publication.

\paragraph{The bill market.} Treasury bills outstanding are of the order of \$6--7\,tn, with daily trading volume
in the hundreds of billions; money market funds and foreign official holders are the largest holders, and primary
dealers intermediate secondary trading subject to balance-sheet constraints that proved binding in March 2020
\citep{duffie2020,vissing2021,he2022}. Two features matter for a stablecoin run. First, bills are the most liquid
instrument in the world, so the price impact of even large flows is measured in basis points, not percent; the
empirical elasticity of three-month bill yields to stablecoin flows is a few basis points per \$3--4\,bn over ten
days, with outflows mattering more than inflows \citep{ahmed2025}. Second, a bill portfolio self-liquidates: a
\$120\,bn ladder of bills with an average maturity of one to two months returns about \$3\,bn of cash every
business day without any sale. An issuer facing redemptions can meet them by not rolling its maturities, which
withdraws demand at auction rather than adding supply in the secondary market. Both channels affect yields; only
the second is a fire sale in the usual sense, and only the second touches dealer balance sheets.

\paragraph{The March 2020 analogue.} Prime money market funds lost about 15\% of assets in two weeks in March
2020, sold commercial paper and bills into dislocated markets, and were backstopped by the Money Market Mutual Fund
Liquidity Facility after ten days \citep{li2021,fed_mmlf2020}. We use this episode to discipline the behaviour of
other bill holders and the design of the backstop lever.

\section{Related literature}\label{sec:lit}

The paper sits at the intersection of three literatures. The first studies stablecoin runs. \citet{gorton2023}
argue that privately issued stablecoins are the wildcat banknotes of our time; \citet{anadu2023} document that
stablecoin flows during 2022--23 resembled money market fund runs and flights to safety; \citet{ma2023} show that
the centralisation of arbitrage, which we model as a capital- and capacity-constrained arbitrage bid, trades off
peg stability against run risk; \citet{lyons2023} study what keeps the peg; \citet{davernas2022} and
\citet{bertsch2023} provide theoretical models of stablecoin fragility; \citet{liao2022} and \citet{barthelemy2023}
analyse the growth of stablecoins and their reserve holdings; \citet{bis2025} frame the policy debate around the
singleness, elasticity and integrity of money. The second literature concerns runs and fire sales more generally.
Our threshold rule is a reduced form of the global-games equilibrium of \citet{goldstein2005}, in the tradition of
\citet{diamond1983}, and our fire-sale cost follows \citet{morris2016}, \citet{shleifer2011} and
\citet{greenwood2015}; the stress-testing use of fire-sale externalities follows \citet{cont2017}. The third
literature concerns the Treasury market: the price-impact anchor comes from \citet{ahmed2025}, the market-structure
and March 2020 evidence from \citet{duffie2020}, \citet{vissing2021}, \citet{he2022} and \citet{fleming2003}, and the
decomposition of impact into permanent and temporary components from \citet{kyle1985} and \citet{almgren2001}, with
the square-root shape documented by \citet{bouchaud2018}. Methodologically we follow the agent-based tradition
advocated by \citet{farmer2009} and use one-at-a-time and global sensitivity designs \citep{saltelli2008,morris1991}.
What is new is the bridge: a calibrated model in which run dynamics, redemption plumbing and bill-market impact are
in one place, so that the size of the spillover, its channel and the leverage of policy can be computed rather than
asserted.

\section{Model}\label{sec:model}

Time is discrete in business days $t=1,\dots,T$ with $T=30$. All quantities are vectorised over Monte Carlo
paths. We describe one issuer; the second issuer (Section~\ref{sec:e4}) is identical with a belief spillover.

\subsection{Holders}

There are $N=10{,}000$ holders of five types $k$: retail, exchanges and market makers, DeFi protocols,
institutional treasuries and arbitrageurs, with shares of supply $(30, 30, 20, 15, 5)\%$ and holdings drawn
lognormally within type. Only institutional holders and arbitrageurs have direct access to the issuer. At the
start of the run the public learns that a share $x$ of reserves is \emph{exposed} (for example, deposits at a
failed bank). A share $r\in[0,1]$ of the exposure is actually lost; the market assigns probability $\psi$ that an
unresolved exposure becomes a loss, so the expected loss is $\ell_0 = x\,[r + (1-r)\psi]$ until a resolution day
$t_R$ on which the public learns the realised loss $xr$. Because the issuer pays at par, the loss falls on whoever
has not left: the expected loss share of a remaining holder is
\begin{equation}
b_t \;=\; \ell_t\, S_0 / S_{t-1},
\end{equation}
where $\ell_t$ is the public expected loss as a share of initial supply and $S_t$ is remaining supply. Holder $i$
observes $b_t$ with a private noise $\varepsilon_i\sim N(0,\sigma^2)$ fixed over the episode, $x_{it}=b_t+\varepsilon_i$,
and forms an exit signal
\begin{equation}
s_{it} \;=\; w_2 x_{it} + w_3 H_t + w_4 \mathbf{1}[\text{suspended}_{t-1}] + \nu_t + w_1 (1-p_{t-1}),
\label{eq:signal}
\end{equation}
where $H_t = \lambda_H H_{t-1} + (\text{exit flow}_t / S_0)$ is a decaying herding state, $\nu_t\sim N(0,\sigma_\nu^2)$
is a common daily sentiment shock and $p_{t-1}$ is yesterday's secondary price. Holder $i$ pays attention on a given
day with a type-specific probability $a_k$ and exits if $s_{it}$ exceeds a threshold $\theta_i\sim N(\bar\theta_k,
\tau_k^2)$. An exited holder re-enters, buying tokens back at the market price, if the \emph{core} signal (the first
four terms of~(\ref{eq:signal}), without the depeg term) has fallen below $\theta_i - h$: once fundamentals are fine,
a discount is a buying opportunity rather than a reason to stay away. Direct-access holders who exit join the
issuer's redemption queue at par; all others sell on the secondary market.

The threshold rule is a reduced form of a global game. In \citet{goldstein2005}, agents with private signals about a
fundamental coordinate on a unique cutoff; here the cutoffs are heterogeneous by type because redemption access,
holding size and costs differ, and we calibrate their level rather than derive it. Appendix~\ref{app:benchmark}
shows that with the herding, depeg and suspension channels switched off and $b_t$ held fixed, the share of supply that
exits converges exactly to the closed-form aggregate $\sum_k \omega_k \Phi\!\big((x-\bar\theta_k)/\sqrt{\sigma^2+\tau_k^2}\big)$,
so the simulation nests a transparent benchmark.

\subsection{Secondary market and arbitrage}

Tokens offered on day $t$ are new sales $V_t$ plus the \emph{overhang} $O_{t-1}$ of tokens held by uninformed dip
buyers who resell. Demand at par from re-entering holders is served first. Arbitrageurs then buy
\begin{equation}
B_t \;=\; \min\!\big(V_t + O_{t-1} - \text{re-entry}_t,\; \kappa S_0,\; \max(0,\, h_a c S_0 - Q_{t-1})\big)
\end{equation}
if the discount covers the redemption fee plus their cost, where $\kappa$ is daily arbitrage capital as a share of
supply, $c$ the issuer's daily processing capacity as a share of supply, $h_a$ the number of days of capacity within
which they require their claim to clear, and $Q$ the redemption queue; a suspended issuer removes the arbitrage bid
entirely. The remainder stays as overhang and sets the liquidity discount. The secondary price is
\begin{equation}
p_t \;=\; \max\Big(1 - \min(b_t, \bar d) - \big[\,c_a \mathbf{1}[\text{offered}>0] + \eta\, O_t/S_0 + \delta\,
\mathbf{1}[\text{suspended}_t]\big],\; p_{\min}\Big),
\end{equation}
a fundamental discount equal to the expected loss share of remaining holders plus an inventory-type liquidity
discount proportional to the stock of unabsorbed tokens. Arbitrage purchases join the redemption queue at par.

\subsection{Issuer}

Reserves are held in cash and deposits, overnight reverse repo, money market fund shares, bills of at most one
month ($\text{tb1}$), bills of one to three months ($\text{tb3}$) and other assets. Impairment removes the realised
loss from the least liquid assets first (other, then cash, then money market fund shares, then repo); bills are never
impaired. Each day a uniform ladder matures: $\text{tb1}$ runs off into cash over 21 business days and $\text{tb3}$
rolls into $\text{tb1}$ over 42. The issuer processes $\min(Q_t, cS_0)$ through a waterfall: cash, repo and money
market fund shares (no price impact), then bills shortest-first at the current price concession, then other assets
at a 30\% haircut; if the waterfall cannot raise the cash the issuer is \emph{suspended} for the day. Liquidity
management is explicit: the issuer keeps cash for the pending claims it can pay within $h_c=5$ days and rolls the
rest of the day's maturities back into bills, so the cash buffer is spent, not rebuilt, during a run. Maturities
kept as cash are bill demand \emph{withdrawn} at auction. Net asset value marks bills at the current concession.
A policy switch lets the issuer redeem at net asset value instead of par, in which case the loss is shared pro
rata and $b_t$ does not grow with redemptions.

\subsection{Bill market}

Let $C_t$ be cumulative net demand withdrawn from the bill market, $C_t = \max(0,\, C_{t-1} + \text{sold}_t +
\omega\,\text{withdrawn}_t - \text{buys}_t)$, where a share $\rho$ of each day's redemption payments returns to
bills through money market funds after a lag $L$ (buys), and $\omega$ weights non-rollover relative to a sale.
Dealer inventory $I_t = (1-\alpha) I_{t-1} + \text{sold}_t - \text{buys}_t$ builds only from outright sales and is
absorbed by other holders at rate $\alpha$. The three-month yield deviation from its no-run path is
\begin{equation}
Y_t \;=\; \underbrace{\lambda_P\, g(C_t)\, m_t}_{\text{persistent}} \;+\; \underbrace{(1-\kappa_T) Y^T_{t-1} +
\lambda_T\, g(\text{sold}_t)\, m_t}_{\text{temporary}}, \qquad m_t = 1 + (I_{t-1}/K)^{\gamma},
\label{eq:impact}
\end{equation}
with $g(q)=q$ for $q\le q_0$ and $g(q)=\sqrt{q_0 q}$ beyond the anchor flow $q_0=\$3.5\,\bn$ (square-root impact,
the baseline), or $g(q)=q$ throughout (linear impact, the upper bound). The coefficients are anchored so that a
\$3.5\,bn outflow raises the three-month yield by about 7\,bp after ten days, the outflow elasticity reported by
\citet{ahmed2025} (to be verified), with buys acting at a third of that strength. The fire-sale loss on bills sold is
face value times duration times the yield deviation, with durations of 0.04 and 0.17 years for the two buckets. A
Fed backstop, when active from day two, removes the dealer-capacity penalty and halves the impact coefficients, a
stylised Money Market Mutual Fund Liquidity Facility.

\subsection{Timing, outputs and experiments}

Within a day: beliefs and signals update; holders exit or re-enter; exits route to the queue or the secondary market;
the secondary market clears; the ladder runs off and the issuer processes the queue through the waterfall; the
issuer manages liquidity; the bill market updates; losses, net asset value and all series are recorded. Each
scenario is run for 1{,}000 Monte Carlo paths in the baseline and 200--300 in the grids; a path draws thresholds,
private signals and common shocks, and a run is defined as cumulative redemptions of at least 20\% of initial
supply by day 30. Five experiments follow the pre-registered design: E1 baseline shocks on Issuer~A, E2 scaling of
issuer size, E3 a policy grid, E4 contagion to Issuer~B, and E5 sensitivity.
